% Proposed captions/notes, compiled only in the draft review documents.
\newcommand{\FrameworkCaption}{Conceptual Research Framework for Multimodal Industrial Monitoring}
\newcommand{\FrameworkNote}{The links organize heterogeneous observations around four
architectural and methodological research subquestions. Acoustic observations belong to
the original scope and literature investigation; predictive audio comparisons were not
completed. Modalities were not jointly available in every study. SQ1--SQ3 are parallel
research dimensions and SQ4 supports their assessment. This framework specifies an
investigation, not a single implemented architecture or an established optimal fusion
system.}
\newcommand{\GapCaption}{Literature Gap Map for Multimodal Industrial and Robotic Monitoring}
\newcommand{\GapNote}{The arrows connect findings under their published conditions to
task-specific unresolved comparisons and five proposed analytical requirements. The
central gap summarizes the comparisons; requirements respond to it. The synthesis is
bounded by the reviewed corpus. R1--R5 do not imply that all historical experiments were
prospectively designed to meet them.}
\newcommand{\GapSources}{%
[1] \citep{Petrich2021Multimodal,Pandiyan2022Deep,Chen2023Multisensor} (physical target
references and correspondence).
[2] \citep{Horwitz2023Back,Roth2022Towards,Nie2023PatchTST} (inspection and transferable
forecasting).
[3] \citep{P003,Hua2026Weld,Thoduka2024Handover,Wang2020What,Ting1999Issues}
(welding/robotics; optimization and stacking foundations).
[4] \citep{P005,P007,Huo2022Multisensor,Altinses2023Fuzzy,Zhang2023QMF,P022,P026}
(physical diagnosis, simulation, transferable quality methods and inspection).
[5] \citep{P004,P017,P008,Bengio2004CVVariance,P026,P033} (methodological evidence and
target-specific inspection).}
\newcommand{\FlowCaption}{Documented Development of the Literature Review Corpus}
\newcommand{\FlowNote}{Counts distinguish screened candidate records, discovery
inclusions, full-text-appraised works and later PDF inputs. The two
abstract/metadata-only inclusions supply no full-text methodological findings and are
outside the initial 37 appraisals. The 94 records comprise 37 initial and 57 additional
work identities, not 94 local PDFs or equivalent industrial fusion experiments. The
dashed abstract/metadata box identifies records outside that historical full-text count.
This is descriptive corpus accounting, not a PRISMA flow.}
\newcommand{\TraceCaption}{Traceability of Proposed Analytical Requirements Across
Research Questions and Empirical Investigations}
\newcommand{\TraceNote}{Coverage identifies comparisons investigated, not requirements
fully satisfied. R1--R5 are proposed synthesis labels, while SQ1--SQ4 retain the
original questions. Methods and checks differ across datasets; the table does not
describe a common prospective protocol. CNN means convolutional neural network; F2 is
the REASSEMBLE logistic stacker. No universal architecture ranking or unchanged-model
cross-domain generalization is established.}
